% ---------------------------------------------------------------------------
% Lecture: Conditional Expectation. Built on one picture: the two-dice table.
% Rows are the events {Y=y}. Homework: quiz03.tex (problems in
% quiz03-problems.tex, numbering as of 2026-09-16).
%
% Every number on these slides was checked with python3 (fractions/sympy);
% scripts: scratchpad/check.py (first pass) and scratchpad/recheck.py
% (revision after critique). Key values, all exact:
%   E[X|Y=y]=y+3.5, E[X]=7, E[X^2|Y=y]=y^2+7y+91/6, E[X^2]=329/6, Var X=35/6
%   E[XY]=329/12=27.4166.., Cov(X,Y)=Cov(X,Z)=35/12, E[X 1_{Y=6}]=57/36=19/12
%   fair die: E=7/2, E[.^2]=91/6, Var=35/12
%   E[Z|X=x]=x/2; two-stage point: E[X|Y=y]=y/2, E[X]=1/4, f_X=-ln x,
%   E[X^2|Y=y]=y^2/3, E[X^2]=1/9, Var X=7/144, E[Y|X=x]=(1-x)/(-ln x)
%   (=0.7213 at x=1/2)
%   mirror (X uniform on [y,1]): E[X|Y=y]=(1+y)/2, f_X=-ln(1-x),
%   E[Y|X=x]=1-x/(-ln(1-x)) (=0.279 at x=1/2)
%   coins: E[H]=7/4, p_H(0)=63/384, p_H(1)=120/384, E[NH]=91/12,
%   Cov(N,H)=35/24, Var N=35/12, slope 1/2
%   N dice: E[S]=49/4, 6+...+6^6=55986 outcomes; die-coin game: E=1
%   random walk: D=xi6+xi7+xi8 takes -3,-1,1,3 w.p. 1/8,3/8,3/8,1/8,
%   E[D]=0, E[D^2]=9/4+3/4=3; E[S_8|S_5]=S_5, E[S_8^2|S_5]=S_5^2+3;
%   given S_5=3, S_8 takes 0,2,4,6 w.p. 1/8,3/8,3/8,1/8 (mean 3, 2nd mom 12)
%   peeking sum: E[S]=6 but mu E[N]=77/12; first-six total 21
%   waiting: HH 6 (P(T=n)=F_{n-1}/2^n), HT 4, HHH 14, 66 42, die 6, E[T|T>3]=9
%   gambler: E[R]=2, E[G]=0
% ---------------------------------------------------------------------------
\newcommand{\Var}{\operatorname{Var}}
\newcommand{\Cov}{\operatorname{Cov}}

% Cell labels of the 6x6 table. #1 = y, #2 = z, #3 = mode:
%   X -> y+z;  G -> y+3.5;  E -> 7;  Z -> z
\newcommand{\dicecell}[3]{%
  \pgfmathtruncatemacro{\s}{#1+#2}\pgfmathsetmacro{\g}{#1+3.5}%
  \if#3X\node at (#2,-#1) {\s};\fi
  \if#3G\node at (#2,-#1) {\g};\fi
  \if#3E\node at (#2,-#1) {7};\fi
  \if#3Z\node at (#2,-#1) {#2};\fi}
% Grid lines on the cell boundaries (cells are centred on integer coordinates).
\newcommand{\dicegrid}{\draw[soft,shift={(0.5,0.5)}] (0,-1) grid[step=1] (6,-7);}

% The 6x6 table. #1 (optional): cell size in cm. #2: row to highlight (0 = none).
% #3: cell label mode as in \dicecell.
\newcommand{\dicetable}[3][0.46]{%
  \begin{tikzpicture}[x=#1cm,y=#1cm,font=\scriptsize]
    \foreach \y in {1,...,6}{
      \ifnum#2=\y \fill[shade] (0.5,-\y+0.5) rectangle (6.5,-\y-0.5);\fi
      \foreach \z in {1,...,6}{\dicecell{\y}{\z}{#3}}}
    \dicegrid
    \foreach \y in {1,...,6}{\node[left,soft] at (0.4,-\y) {$y{=}\y$};}
    \foreach \z in {1,...,6}{\node[soft] at (\z,-0.15) {$\z$};}
    \node[left,soft] at (0.4,-0.15) {$z{=}$};
  \end{tikzpicture}}

% The table without labels. #1 (optional): cell size. #2: mode.
\newcommand{\minitable}[2][0.42]{%
  \begin{tikzpicture}[x=#1cm,y=#1cm,font=\scriptsize]
    \foreach \y in {1,...,6}{\foreach \z in {1,...,6}{\dicecell{\y}{\z}{#2}}}
    \dicegrid
  \end{tikzpicture}}

% The X table with each row's average written to its right.
\newcommand{\rowavgtable}[1][0.4]{%
  \begin{tikzpicture}[x=#1cm,y=#1cm,font=\scriptsize]
    \foreach \y in {1,...,6}{
      \pgfmathsetmacro{\g}{\y+3.5}
      \foreach \z in {1,...,6}{\dicecell{\y}{\z}{X}}
      \node[right,accent] at (6.6,-\y) {$\to\g$};
    }
    \dicegrid
    \foreach \y in {1,...,6}{\node[left,soft] at (0.4,-\y) {$y{=}\y$};}
  \end{tikzpicture}}

\begin{frame}{Conditional Expectation}
  \small
  One picture carries the whole lecture: the table of outcomes of two dice, $Y$ the first, $Z$ the second, $X=Y+Z$. The rows of the table are the events $\{Y=y\}$.
  \bigskip
  \renewcommand{\arraystretch}{1.8}
  \begin{tabular}{@{}r@{\ \ }>{\raggedright\arraybackslash}p{4.2cm}>{\raggedright\arraybackslash}p{8.6cm}@{}}
    1. & \textbf{Given a value:} $\E[X\mid Y=y]$ & The average of $X$ along row $y$. A number for each $y$, hence a function of $y$. \\
    2. & \textbf{Law of total expectation} & The average of the row averages, weighted by row probabilities, is the overall average. \\
    3. & \textbf{Given the variable:} $\E[X\mid Y]$ & Paint every row with its own average: a new random variable. \\
    4. & \textbf{Properties} & Tower, taking out what is known, independence. Each is a one-glance fact about rows. \\
    5. & \textbf{Two applications} & Random sums, and waiting times: hard head-on, one line by conditioning. \\
  \end{tabular}
\end{frame}

% ---------------------------------------------------------------------------
\begin{frame}{Three Problems We Cannot Yet Do}
  \small
  Each asks for an expected value. With the tools so far, the only route is to find the distribution of the quantity first.
  \begin{block}{Problem A: a random number of coins}
    Roll a die, getting $N$. Toss $N$ fair coins; $H$ is the number of heads. Find $\E[H]$.\\
    Head-on: $p_H(h)=\sum_{n=\max(h,1)}^{6}\tfrac16\binom nh2^{-n}$ for $h=0,\dots,6$, seven sums, then $\sum_h h\,p_H(h)$.
  \end{block}
  \begin{block}{Problem B: an insurer's payout}
    An insurer receives $N$ claims in a year, $\E[N]=20$, each of mean size $\pounds1500$, sizes independent of the count. Find the expected total payout $S=X_1+\dots+X_N$.\\
    Head-on: the distribution of $S$ mixes, over $n$, the $n$-fold convolution of the claim-size distribution, and neither distribution was given.
  \end{block}
  \begin{block}{Problem C: waiting for HH}
    Toss a fair coin until two heads in a row appear; $T$ is the number of tosses used. Find $\E[T]$.\\
    Head-on: $\Prob(T=n)=F_{n-1}/2^n$ with $F_1,F_2,F_3,\dots=1,1,2,3,5,\dots$ the Fibonacci numbers, then $\sum_n n\,\Prob(T=n)$.
  \end{block}
  All three are solved by conditioning on the first stage: slides~\ref{coins}, \ref{insurer}, \ref{hh}.
\end{frame}

% ===========================================================================
\begin{frame}{The Table}
  \small
  \begin{columns}[T]
    \begin{column}{0.52\textwidth}
      Two fair dice: $(\Omega,\F,\Prob)$ with $\Omega=\{1,\dots,6\}^2$, $\F$ all subsets of $\Omega$, and $\Prob(\{\omega\})=\tfrac1{36}$ for each outcome $\omega=(y,z)$.
      \begin{itemize}
        \item $Y\colon\Omega\to\{1,\dots,6\}$, $Y(\omega)=y$, the first die: the \textbf{row}.
        \item $Z\colon\Omega\to\{1,\dots,6\}$, $Z(\omega)=z$, the second die: the \textbf{column}.
        \item $X\colon\Omega\to\{2,\dots,12\}$, $X(\omega)=y+z$, the sum: the \textbf{number in the cell}.
      \end{itemize}
      \smallskip
      Shorthand, used throughout: $\{Y=y\}:=\ev{Y(\omega)=y}$ and $\Prob(Y=y):=\Prob(\{Y=y\})$; likewise $\{X\in A\}$, $\Prob(X\le x)$, $\Prob(X=x,Y=y)$.\par
      \smallskip
      Every random variable on this $\Omega$ writes a number in each of the $36$ cells; every event is a set of cells, with probability the number of cells over $36$.
    \end{column}
    \begin{column}{0.44\textwidth}
      \centering
      $X(\omega)=y+z$\par\medskip
      \dicetable[0.46]{0}{X}
    \end{column}
  \end{columns}
\end{frame}

% ---------------------------------------------------------------------------
\begin{frame}{Conditioning on $Y=y$ Is Restricting to a Row}
  \small
  \label{condpmf}%
  \begin{columns}[T]
    \begin{column}{0.54\textwidth}
      The event $\{Y=4\}$ is the fourth row: six cells, probability $\tfrac6{36}=\tfrac16$.
      Given $Y=4$, the outcome lies in that row, and each of its six cells has conditional probability
      \[
        \Prob(\{\omega\}\mid Y=4)=\frac{\Prob(\{\omega\}\cap\{Y=4\})}{\Prob(Y=4)}=\frac{1/36}{1/6}=\frac16 .
      \]
      So, given $Y=4$, $X$ takes each of the values $5,6,\dots,10$ with probability $\tfrac16$: the distribution of the numbers in row $4$.
    \end{column}
    \begin{column}{0.42\textwidth}
      \centering
      row $4$ is the event $\{Y=4\}$\par\medskip
      \dicetable[0.4]{4}{X}
    \end{column}
  \end{columns}
  \begin{block}{Definition (conditional pmf)}
    Let $S_Y=\{y : p_Y(y)>0\}$. The \textbf{conditional pmf of $X$ given $Y$} is
    \[
      p_{X\mid Y}\colon\R\times S_Y\to[0,1],\qquad p_{X\mid Y}(x\mid y)=\Prob(X=x\mid Y=y)=\frac{p_{X,Y}(x,y)}{p_Y(y)} ;
    \]
    the first argument $x$ is a value of $X$, the second argument $y$ names the row.
  \end{block}
  \words{Restrict to row $y$ and rescale its probabilities so that they total $1$.}
\end{frame}

% ---------------------------------------------------------------------------
\begin{frame}{Conditional Expectation Given $Y=y$: The Row Average}
  \small
  \begin{block}{Definition (discrete case)}
    Let $S_Y=\{y : p_Y(y)>0\}$. The \textbf{conditional expectation of $X$ given $Y=y$} is the function
    \[
      g\colon S_Y\to\R,\qquad g(y)=\E[X\mid Y=y]=\sum_x x\,p_{X\mid Y}(x\mid y).
    \]
  \end{block}
  \words{The ordinary expected value, computed with the row's conditional distribution in place of the whole table's.}
  \begin{block}{What each symbol is, and what the object is not}
    \begin{itemize}
      \item $x$: runs over the values of $X$, summed away. \ $y$: held fixed, names the row. \ Result: depends on $y$ alone.
      \item $g\colon S_Y\to\R$: a function on the values of $Y$, not on $\Omega$. Not a random variable.
      \item $y\notin S_Y$: undefined. A row of probability $0$ has no conditional distribution.
      \item $\E[X\,\ind_{\{Y=y\}}]=\sum_x x\,p_{X,Y}(x,y)=\E[X\mid Y=y]\,\Prob(Y=y)$, with $\ind_{\{Y=y\}}\colon\Omega\to\{0,1\}$ the indicator of the row: the row's total, not its average.
    \end{itemize}
  \end{block}
\end{frame}

% ---------------------------------------------------------------------------
\begin{frame}{Two Dice: The Row Average as a Function of $y$}
  \small
  \label{rowavg}%
  \begin{columns}[T]
    \begin{column}{0.56\textwidth}
      \begin{exampleblock}{$\E[X\mid Y=y]$ for the sum of two dice}
        Row $y$ contains $y+1,\dots,y+6$, each with conditional probability $\tfrac16$:
        \[
          \E[X\mid Y=y]=\tfrac16\big((y+1)+\dots+(y+6)\big)=y+3.5 .
        \]
        As a function on $S_Y=\{1,\dots,6\}$:\par\smallskip
        \centering
        \setlength{\tabcolsep}{4pt}
        \begin{tabular}{@{}c|cccccc@{}}
          $y$ & 1 & 2 & 3 & 4 & 5 & 6 \\ \hline
          $g(y)=\E[X\mid Y=y]$ & 4.5 & 5.5 & 6.5 & 7.5 & 8.5 & 9.5
        \end{tabular}
      \end{exampleblock}
      \begin{itemize}
        \item $g\colon\{1,\dots,6\}\to\R$, $g(y)=y+3.5$. Six numbers, one per row.
        \item $g(y)-y=3.5=\E[Z]$ for every $y$: on row $y$ the first die is fixed at $y$ and all the remaining randomness is the second die. Slide~\ref{constrow} makes this a method.
      \end{itemize}
    \end{column}
    \begin{column}{0.40\textwidth}
      \centering
      each row, then its average\par\medskip
      \rowavgtable
    \end{column}
  \end{columns}
\end{frame}

% ---------------------------------------------------------------------------
\begin{frame}{The Continuous Case: Rows Become Slices}
  \small
  \begin{block}{Definition (continuous case)}
    Let $X,Y$ have joint pdf $f_{X,Y}\colon\R^2\to[0,\infty)$, let $S_Y=\{y : f_Y(y)>0\}$, and define the \textbf{conditional pdf of $X$ given $Y$}
    \[
      f_{X\mid Y}\colon\R\times S_Y\to[0,\infty),\qquad f_{X\mid Y}(x\mid y)=\frac{f_{X,Y}(x,y)}{f_Y(y)},
    \]
    $x$ the value of $X$, $y$ the slice. The \textbf{conditional expectation of $X$ given $Y=y$} is the function
    \[
      g\colon S_Y\to\R,\qquad g(y)=\E[X\mid Y=y]=\int_{-\infty}^{\infty}x\,f_{X\mid Y}(x\mid y)\dd x .
    \]
  \end{block}
  \words{The row $\{Y=y\}$ is now a horizontal line through the plane, and the slice of $f_{X,Y}$ along it, rescaled to area $1$, plays the part of the row's conditional pmf.}
  \begin{itemize}
    \item For fixed $y\in S_Y$, $f_{X\mid Y}(\,\cdot\mid y)$ is a pdf in $x$: $\int f_{X\mid Y}(x\mid y)\dd x=f_Y(y)/f_Y(y)=1$.
    \item $\Prob(Y=y)=0$, so $\Prob(\,\cdot\mid Y=y)$ cannot be defined by dividing probabilities; the ratio of densities is the limit of conditioning on $\{y\le Y\le y+\delta\}$ as $\delta\to0$.
    \item Otherwise nothing changes: fix $y$, take the slice, rescale, average.
  \end{itemize}
\end{frame}

% ---------------------------------------------------------------------------
\begin{frame}{The Two-Stage Point: Horizontal and Vertical Slices}
  \small
  \begin{columns}[T]
    \begin{column}{0.66\textwidth}
      $Y$ uniform on $[0,1]$; given $Y=y$, $X$ uniform on $[0,y]$. So $f_Y=1$ on $(0,1]$, $f_{X\mid Y}(x\mid y)=\tfrac1y$ on $[0,y]$, and $f_{X,Y}(x,y)=f_{X\mid Y}(x\mid y)f_Y(y)=\tfrac1y$ on the triangle $0<x<y<1$.
      \begin{exampleblock}{Given $Y=y$: a horizontal slice}
        $S_Y=(0,1]$, the slice is uniform on $[0,y]$, and
        $\E[X\mid Y=y]=\int_0^y x\,\tfrac1y\dd x=\tfrac y2$.
      \end{exampleblock}
      \begin{exampleblock}{Given $X=x$: a vertical slice}
        $f_X(x)=\int_x^1\tfrac1y\dd y=-\ln x$ on $S_X=(0,1)$, so $f_{Y\mid X}(y\mid x)=\dfrac{1/y}{-\ln x}$ for $x<y<1$ and
        $\E[Y\mid X=x]=\displaystyle\int_x^1\frac{y}{y}\,\frac{\dd y}{-\ln x}=\frac{1-x}{-\ln x}$.
      \end{exampleblock}
      \begin{itemize}
        \item The vertical slice is not uniform, $f_{Y\mid X}(y\mid x)\propto\tfrac1y$ favours small $y$, so $\E[Y\mid X=x]$ lies below the slice's midpoint $\tfrac{1+x}2$: at $x=\tfrac12$, $0.7213$ against $0.75$.
      \end{itemize}
    \end{column}
    \begin{column}{0.32\textwidth}
      \centering
      \begin{tikzpicture}[scale=2.6,font=\scriptsize]
        \fill[shade!50] (0,0) -- (1,1) -- (0,1) -- cycle;
        \draw[->] (0,0) -- (1.15,0) node[right] {$x$};
        \draw[->] (0,0) -- (0,1.15) node[above] {$y$};
        \draw[soft] (0,0) -- (1,1) node[right,pos=0.3,black] {\ $x=y$};
        \draw[accent,very thick] (0,0.55) -- (0.55,0.55);
        \node[accent,left] at (0,0.55) {$y$};
        \fill[accent] (0.275,0.55) circle (0.6pt) node[above,accent] {$\tfrac y2$};
        \draw[leaf,very thick] (0.7,0.7) -- (0.7,1);
        \node[leaf,below] at (0.7,0) {$x$};
        \draw[soft,dashed] (0.7,0) -- (0.7,0.7);
        \fill[leaf] (0.7,0.84) circle (0.6pt) node[right,leaf] {$\tfrac{1-x}{-\ln x}$};
        \node[below] at (1,0) {$1$};
        \node[left] at (0,1) {$1$};
      \end{tikzpicture}
      \par\medskip
      {\footnotesize Blue: $\{Y=y\}$, uniform on $[0,y]$. Green: $\{X=x\}$, density $\propto\tfrac1y$ on $[x,1]$.}
    \end{column}
  \end{columns}
\end{frame}

% ---------------------------------------------------------------------------
\begin{frame}{Functions of $X$, and What Is Constant on a Row}
  \small
  \label{constrow}%
  The same recipe applies to any random variable, and in particular to any function of $X$: for $W\colon\Omega\to\R$ and $h\colon\R\to\R$,
  \[
    \E[W\mid Y=y]=\sum_w w\,p_{W\mid Y}(w\mid y),\qquad
    \E[h(X)\mid Y=y]=\sum_x h(x)\,p_{X\mid Y}(x\mid y),
  \]
  with $p_{W\mid Y}$ defined as on slide~\ref{condpmf} with $W$ in place of $X$; the second formula is the rule $\E[h(X)]=\sum_x h(x)p_X(x)$ applied to the row's pmf.
  \begin{block}{Observation: on row $y$, $Y$ is the constant $y$}
    Anything that depends on the outcome only through $Y$ is constant on the row. Given $Y=y$, $X=Y+Z$ is $y+Z$: the row's randomness is all in $Z$.
  \end{block}
  \begin{exampleblock}{Two dice: $\E[X^2\mid Y=y]$ without listing the row}
    On row $y$, $X^2=(y+Z)^2=y^2+2yZ+Z^2$ with $y$ a constant and $Z$ uniform on $\{1,\dots,6\}$, so
    \[
      \E[X^2\mid Y=y]=y^2+2y\,\E[Z]+\E[Z^2]=y^2+7y+\tfrac{91}{6} .
    \]
    Head-on you would square six numbers and average, for each of six rows.
  \end{exampleblock}
  \begin{itemize}
    \item The observation is the seed of the property ``taking out what is known'' (slide~\ref{takeout}).
  \end{itemize}
\end{frame}

% ===========================================================================
\begin{frame}{Total Expectation: The Average of the Row Averages}
  \small
  \begin{columns}[T]
    \begin{column}{0.56\textwidth}
      The rows partition the table. The overall average of the $36$ cells can be computed row by row:
      \[
        \E[X]=\sum_{y=1}^6 \underbrace{\tfrac16}_{\Prob(Y=y)}\cdot\underbrace{(y+3.5)}_{\text{row average}}=7 .
      \]
      \begin{block}{Theorem (law of total expectation)}
        \begin{align*}
          \text{discrete:}&\quad \E[X]=\sum_{y\in S_Y}\E[X\mid Y=y]\,p_Y(y),\\
          \text{continuous:}&\quad \E[X]=\int_{S_Y}\E[X\mid Y=y]\,f_Y(y)\dd y .
        \end{align*}
      \end{block}
      \words{Average each row, then average the row averages weighted by the row probabilities; because the rows cover the table exactly once, this is the overall average.}
    \end{column}
    \begin{column}{0.40\textwidth}
      \centering
      \begin{tikzpicture}[x=0.4cm,y=0.4cm,font=\scriptsize]
        \foreach \y in {1,...,6}{
          \pgfmathsetmacro{\g}{\y+3.5}
          \foreach \z in {1,...,6}{\dicecell{\y}{\z}{X}}
          \node[right,accent] at (6.6,-\y) {$\to\g$};
        }
        \dicegrid
        \node[accent,font=\small] at (7.9,-7.3) {$\downarrow$};
        \node[accent] at (7.9,-8.2) {$\tfrac16\sum=7$};
      \end{tikzpicture}
    \end{column}
  \end{columns}
\end{frame}

% ---------------------------------------------------------------------------
\begin{frame}{Total Expectation: Proof}
  \small
  \begin{block}{Proof (discrete case)}
    \begin{align*}
      \sum_y \E[X\mid Y=y]\,p_Y(y)
      &=\sum_y\Big(\sum_x x\,p_{X\mid Y}(x\mid y)\Big)p_Y(y) && \text{\footnotesize definition}\\
      &=\sum_y\sum_x x\,p_{X,Y}(x,y) && \text{\footnotesize undo the rescaling}\\
      &=\sum_x x\sum_y p_{X,Y}(x,y)=\sum_x x\,p_X(x)=\E[X]. && \text{\footnotesize swap sums; marginal}\quad\square
    \end{align*}
  \end{block}
  \words{Line two undoes the rescaling that defined the conditional pmf, and line three sums the cells column-first instead of row-first.}
  \begin{itemize}
    \item The swap of the two sums is free on the finite table; in general it needs $\E|X|<\infty$.
    \item Continuous case: identical, with $f$ for $p$ and $\int$ for $\sum$:
      $\int_{S_Y}\Big(\int x\,f_{X\mid Y}(x\mid y)\dd x\Big)f_Y(y)\dd y=\iint x\,f_{X,Y}(x,y)\dd x\dd y=\int x\,f_X(x)\dd x$.
    \item Any $W\colon\Omega\to\R$ in place of $X$: $\E[W]=\sum_y\E[W\mid Y=y]\,p_Y(y)$, in particular for $W=h(X)$.
  \end{itemize}
\end{frame}

% ---------------------------------------------------------------------------
\begin{frame}{Total Expectation at Work (1): Problem A, a Random Number of Coins}
  \small
  \label{coins}%
  Roll a die, getting $N\colon\Omega\to\{1,\dots,6\}$. Toss $N$ fair coins, independent of the die; $H\colon\Omega\to\{0,\dots,6\}$ is the number of heads. What is $\E[H]$?
  \begin{block}{By conditioning}
    Given $N=n$, $H$ is binomial with $n$ tosses (the coins are independent of the die, so the row $\{N=n\}$ does not change their distribution), so $\E[H\mid N=n]=\tfrac n2$. Then
    \[
      \E[H]=\sum_{n=1}^6\tfrac n2\cdot\tfrac16=\tfrac12\,\E[N]=1.75 .
    \]
  \end{block}
  \begin{block}{Head-on, for comparison}
    $p_H(h)=\sum_{n=\max(h,1)}^{6}\tfrac16\binom nh 2^{-n}$ for $h=0,\dots,6$: seven sums (e.g.\ $p_H(0)=\tfrac{63}{384}$, $p_H(1)=\tfrac{120}{384}$), then $\sum_h h\,p_H(h)$.
  \end{block}
  \begin{itemize}
    \item Given the row $N=n$ the problem is a textbook one. The distribution of $H$, a mixture of six binomials, was never needed.
  \end{itemize}
\end{frame}

% ---------------------------------------------------------------------------
\begin{frame}{Total Expectation at Work (2): The Two-Stage Point, and Variance}
  \small
  \label{variance}%
  $Y$ uniform on $[0,1]$; given $Y=y$, $X$ uniform on $[0,y]$. What is $\E[X]$?
  \begin{block}{By conditioning}
    $\E[X\mid Y=y]=\tfrac y2$ and $f_Y=1$ on $(0,1]$, so
    $\E[X]=\int_0^1\tfrac y2\cdot1\dd y=\tfrac14$.
  \end{block}
  \begin{block}{Head-on, for comparison}
    First find the pdf of $X$ by integrating out $y$: $f_X(x)=\int_x^1 f_{X,Y}(x,y)\dd y=\int_x^1\tfrac{\dd y}{y}=-\ln x$ for $0<x<1$.
    Then $\E[X]=\int_0^1 x(-\ln x)\dd x$, which needs integration by parts, and gives $\tfrac14$.
  \end{block}
  \begin{block}{Definition (variance)}
    For $X\colon\Omega\to\R$ with $\E[X^2]<\infty$: \quad $\Var(X)=\E\big[(X-\E[X])^2\big]=\E[X^2]-\E[X]^2$.
  \end{block}
  \words{The average squared distance of $X$ from its mean; computing it needs $\E[X^2]$, which total expectation supplies row by row.}
  \begin{itemize}
    \item Here $\E[X^2\mid Y=y]=\tfrac{y^2}3$, so $\E[X^2]=\int_0^1\tfrac{y^2}3\dd y=\tfrac19$ and $\Var(X)=\tfrac19-\tfrac1{16}=\tfrac7{144}$. Neither computation needed $f_X$.
  \end{itemize}
\end{frame}

% ===========================================================================
\begin{frame}{$\E[X\mid Y]$: Paint Every Row With Its Average}
  \small
  \begin{columns}[T]
    \begin{column}{0.54\textwidth}
      Slide~\ref{rowavg} gave a function on the values of $Y$: $g\colon S_Y\to\R$, $g(y)=\E[X\mid Y=y]$.
      \medskip

      Now write $g(y)$ into \emph{every cell of row $y$}. The result is a new number in each cell of the table, that is, a new random variable.
      \begin{block}{Definition}
        The \textbf{conditional expectation of $X$ given $Y$} is the random variable
        \begin{align*}
          \E[X\mid Y]\colon\Omega&\to\R,\\
          \E[X\mid Y](\omega)&=g\big(Y(\omega)\big)=\E\big[X\mid Y=Y(\omega)\big].
        \end{align*}
      \end{block}
      \words{At outcome $\omega$, look up which row $\omega$ is in and report that row's average.}
      \begin{itemize}
        \item $g(Y)$ is defined on the event $\{Y\in S_Y\}$, which has probability $1$; set it to $0$ elsewhere. On the two-dice $\Omega$ every outcome has positive probability, so there is nothing to fix.
      \end{itemize}
    \end{column}
    \begin{column}{0.42\textwidth}
      \centering
      $\E[X\mid Y](\omega)=Y(\omega)+3.5$\par\medskip
      \dicetable[0.5]{0}{G}
    \end{column}
  \end{columns}
\end{frame}

% ---------------------------------------------------------------------------
\begin{frame}{Three Tables, Three Levels of Knowledge}
  \small
  \begin{center}
    \begin{tabular}{@{}c@{\quad}c@{\quad}c@{}}
      $X$ & $\E[X\mid Y]$ & $\E[X]$ \\[2pt]
      \dicetable[0.48]{0}{X} & \dicetable[0.48]{0}{G} & \dicetable[0.48]{0}{E} \\[4pt]
      you see the outcome & you see the row only & you see nothing
    \end{tabular}
  \end{center}
  \words{Each table is the previous one averaged over what you do not know.}
  \begin{itemize}
    \item $\E[X\mid Y=y]$: a number for each $y$. \ $\E[X\mid Y]=g(Y)$: a random variable, constant on each row $\{Y=y\}$, where its value is $\E[X\mid Y=y]$.
  \end{itemize}
\end{frame}

% ---------------------------------------------------------------------------
\begin{frame}{$\E[X\mid Y]$: Extreme Cases and the Continuous Case}
  \small
  \begin{itemize}
    \item \textbf{$\E[X\mid X]=X$.} Condition on $X$ itself: each ``row'' $\{X=x\}$ contains only cells where $X=x$, so its average is $x$. Full knowledge changes nothing.
    \item \textbf{$\E[X\mid Y]=\E[X]$ if $Y$ is constant.} One row, the whole table. No knowledge.
    \item \textbf{$\E[Y\mid Y]=Y$}, and more generally $\E[h(Y)\mid Y]=h(Y)$ for $h\colon\R\to\R$: on row $y$, $h(Y)$ is the constant $h(y)$, whose average is $h(y)$.
  \end{itemize}
  \begin{block}{Continuous case}
    The same definition, $\E[X\mid Y]=g(Y)$ with $g\colon S_Y\to\R$, $g(y)=\int x\,f_{X\mid Y}(x\mid y)\dd x$. The ``rows'' $\{Y=y\}$ each have probability $0$, but $g(Y)\colon\Omega\to\R$ is an ordinary random variable: a function of $Y$.
  \end{block}
  \begin{exampleblock}{Examples}
    Two-stage point: $g(y)=y/2$, so $\E[X\mid Y]=Y/2$. At an outcome where $Y=0.8$, the average for $X$ is $0.4$, whatever $X$ actually is. Likewise $\E[Y\mid X]=(1-X)/(-\ln X)$.\\
    Random coins: $\E[H\mid N]=N/2$, a random variable with values $\tfrac12,1,\dots,3$, each with probability $\tfrac16$.
  \end{exampleblock}
\end{frame}

% ===========================================================================
\begin{frame}{Property 1: Tower. The Painted Table Has the Same Average}
  \small
  \begin{columns}[T]
    \begin{column}{0.54\textwidth}
      Painting a row with its average does not change the row's total. So the painted table has the same overall average as the original:
      \begin{block}{Theorem (tower property)}
        $\E\big[\E[X\mid Y]\big]=\E[X]$.
      \end{block}
      \begin{block}{Proof}
        $\E[X\mid Y]=g(Y)$, and the expectation of a function of $Y$ is
        $\E[g(Y)]=\sum_y g(y)\,p_Y(y)=\sum_y\E[X\mid Y=y]\,p_Y(y)=\E[X]$ by total expectation. \hfill$\square$
      \end{block}
      \words{It is the law of total expectation with the $y$'s hidden inside a random variable.}
      \begin{itemize}
        \item That makes it composable: $\E[X\mid Y]$ is itself a random variable, and can be attacked by conditioning on something else.
      \end{itemize}
    \end{column}
    \begin{column}{0.44\textwidth}
      \centering
      \begin{tabular}{@{}c@{\ }c@{\ }c@{}}
        \minitable[0.44]{X} & {\color{accent}$\to$} & \minitable[0.44]{G} \\[2pt]
        average $7$ & & average $7$
      \end{tabular}
    \end{column}
  \end{columns}
\end{frame}

% ---------------------------------------------------------------------------
\begin{frame}{Property 1: Why the Tower Matters}
  \small
  To compute $\E[X]$, compute $\E[X\mid Y]$ for a well-chosen $Y$ and average that instead.
  \begin{center}
    \renewcommand{\arraystretch}{1.4}
    \begin{tabular}{@{}lll@{}}
      \toprule
      problem & $\E[X\mid Y]$, a formula in $Y$ & then $\E[X]=\E[\E[X\mid Y]]$ \\
      \midrule
      coins, $N$ from a die & $\E[H\mid N]=N/2$ & $\E[N]/2=1.75$ \\
      two-stage point & $\E[X\mid Y]=Y/2$ & $\E[Y]/2=1/4$ \\
      random sum $S=X_1+\dots+X_N$ & $\E[S\mid N]=\mu N$ & $\mu\,\E[N]$ \\
      waiting time $T$ for HH & $\E[T\mid Y]$ for the opening $Y$ & an equation for $\E[T]$ \\
      \bottomrule
    \end{tabular}
  \end{center}
  \begin{itemize}
    \item In each row $X$ has an unpleasant distribution; $\E[X\mid Y]$ is a one-line formula in $Y$.
    \item Choosing $Y$ is the whole art: the row problem must be textbook, and $p_Y$ or $f_Y$ must be known.
    \item The last two rows are the applications at the end of the lecture.
  \end{itemize}
\end{frame}

% ---------------------------------------------------------------------------
\begin{frame}{Property 2: A Problem. Are $N$ and $H$ Correlated?}
  \small
  Problem A again: $N$ the die, $H$ the number of heads among $N$ coins. More coins should mean more heads. How strongly are $N$ and $H$ linked?
  \begin{block}{Definition (covariance)}
    For $X,Y\colon\Omega\to\R$ with finite second moments: \quad $\Cov(X,Y)=\E\big[(X-\E[X])(Y-\E[Y])\big]=\E[XY]-\E[X]\,\E[Y]$; \quad $\Cov(X,X)=\Var(X)$.
  \end{block}
  \words{Positive when $X$ above its mean tends to come with $Y$ above its mean; zero for independent $X,Y$.}
  \begin{block}{Head-on}
    $\E[NH]=\sum_{n=1}^6\sum_{h=0}^n n\,h\,p_{N,H}(n,h)$ with $p_{N,H}(n,h)=\tfrac16\binom nh2^{-n}$: $21$ nonzero terms of the joint pmf, each a binomial coefficient over a power of two.
  \end{block}
  \begin{block}{What conditioning offers}
    Tower: $\E[NH]=\E\big[\E[NH\mid N]\big]$. On the row $\{N=n\}$ the factor $N$ is the number $n$. If a constant factor can be pulled out of a row average, then
    $\E[NH\mid N=n]=n\,\E[H\mid N=n]=n\cdot\tfrac n2$,
    and $\E[NH]$ is an ordinary expectation of a function of $N$. Slide~\ref{constrow} made this move one row at a time; the next slide states it once, for all rows.
  \end{block}
\end{frame}

% ---------------------------------------------------------------------------
\begin{frame}{Property 2: Taking Out What Is Known}
  \small
  \label{takeout}%
  On row $y$, anything that is a function of $Y$ is the constant $h(y)$, and a constant factor comes out of an average.
  \begin{block}{Theorem}
    For any $h\colon\R\to\R$ with $\E|h(Y)X|<\infty$: \quad $\E[h(Y)\,X\mid Y]=h(Y)\,\E[X\mid Y]$. \quad In particular $\E[h(Y)\mid Y]=h(Y)$.
  \end{block}
  \begin{block}{Proof (discrete case)}
    Fix $y\in S_Y$. $\E[h(Y)X\mid Y=y]=\sum_x h(y)\,x\,p_{X\mid Y}(x\mid y)=h(y)\sum_x x\,p_{X\mid Y}(x\mid y)=h(y)\,\E[X\mid Y=y]$.
    This holds for every $y$; substitute $y=Y(\omega)$. Continuous case identical, with $f$ for $p$ and $\int$ for $\sum$. \hfill$\square$
  \end{block}
  \words{Whatever is already determined by the row is a constant on the row and can be taken out of the row average.}
  \begin{block}{Linearity}
    $\E[aX+bX'\mid Y]=a\,\E[X\mid Y]+b\,\E[X'\mid Y]$ for constants $a,b$, for the same reason ordinary expectation is linear: row by row.
    With the theorem: $\E[h(Y)+X\mid Y]=h(Y)+\E[X\mid Y]$.
  \end{block}
\end{frame}

% ---------------------------------------------------------------------------
\begin{frame}{Property 2: Back to $\Cov(N,H)$, and $\E[XY]$ for Two Dice}
  \small
  \begin{exampleblock}{$\Cov(N,H)$ for Problem A}
    \[
      \E[NH]=\E\big[\E[NH\mid N]\big]=\E\big[N\,\E[H\mid N]\big]=\E\big[N\cdot\tfrac N2\big]=\tfrac12\E[N^2]=\tfrac12\cdot\tfrac{91}6=\tfrac{91}{12},
    \]
    so $\Cov(N,H)=\tfrac{91}{12}-\tfrac72\cdot\tfrac74=\tfrac{35}{24}\approx1.46$.
    With $\Var(N)=\E[N^2]-\E[N]^2=\tfrac{91}6-\tfrac{49}4=\tfrac{35}{12}$: $\Cov(N,H)/\Var(N)=\tfrac12$, the slope of $\E[H\mid N]=\tfrac12N$.
  \end{exampleblock}
  \begin{exampleblock}{$\E[XY]$ for two dice, without listing $36$ products}
    \[
      \E[XY]=\E\big[\E[XY\mid Y]\big]=\E\big[Y\,\E[X\mid Y]\big]=\E[Y(Y+3.5)]=\E[Y^2]+3.5\,\E[Y]=\tfrac{91}{6}+\tfrac{49}{4}=\tfrac{329}{12},
    \]
    about $27.42$, so $\Cov(X,Y)=\tfrac{329}{12}-7\cdot3.5=\tfrac{35}{12}$.
  \end{exampleblock}
  \begin{itemize}
    \item Both computations are the same three steps: tower; take out the factor that is a function of the conditioning variable; an ordinary expectation of a function of that variable.
    \item The $21$-term double sum of the previous slide was never written down.
  \end{itemize}
\end{frame}

% ---------------------------------------------------------------------------
\begin{frame}{Property 3: Independence}
  \small
  \begin{columns}[T]
    \begin{column}{0.64\textwidth}
      \begin{block}{Definition (independence)}
        $X,Y\colon\Omega\to\R$ are \textbf{independent}, written $X\indep Y$, if
        \[
          \Prob(X\in A,\,Y\in B)=\Prob(X\in A)\,\Prob(Y\in B)\quad\text{for all }A,B\subseteq\R .
        \]
        $\xi_1,\dots,\xi_n\colon\Omega\to\R$ are \textbf{independent} if
        $\Prob(\xi_1\in A_1,\dots,\xi_n\in A_n)=\prod_{i=1}^n\Prob(\xi_i\in A_i)$ for all $A_i\subseteq\R$; a sequence, if every finite part is.
      \end{block}
      \words{The row you are in tells you nothing about which column you are in.}
      \begin{itemize}
        \item Two dice: $p_{Y,Z}(y,z)=\tfrac1{36}=\tfrac16\cdot\tfrac16=p_Y(y)\,p_Z(z)$; summing over the cells of $\{Y\in A\}\cap\{Z\in B\}$ gives the definition for all $A,B$. So $Y\indep Z$.
      \end{itemize}
    \end{column}
    \begin{column}{0.34\textwidth}
      \centering
      $Z(\omega)=z$: every row alike\par\medskip
      \dicetable[0.4]{0}{Z}
    \end{column}
  \end{columns}
  \begin{block}{Fact (stated without proof)}
    Functions of disjoint groups of independent variables are independent: $\xi_1,\xi_2,\dots$ independent implies $f(\xi_1,\dots,\xi_5)\indep g(\xi_6,\xi_7,\xi_8)$ for all $f\colon\R^5\to\R$, $g\colon\R^3\to\R$.
  \end{block}
\end{frame}

% ---------------------------------------------------------------------------
\begin{frame}{Property 3: A Problem. A Random Walk, Seen From Halfway}
  \small
  \label{rw}%
  Let $\xi_1,\xi_2,\dots\colon\Omega\to\{-1,1\}$ be independent with $\Prob(\xi_i=1)=\Prob(\xi_i=-1)=\tfrac12$, and $S_n=\xi_1+\dots+\xi_n$ the position after $n$ steps. You watch the first five steps and see $S_5$. What do you expect $S_8$ to be, and $S_8^2$? That is: find $\E[S_8\mid S_5]$ and $\E[S_8^2\mid S_5]$.
  \begin{block}{Head-on}
    Row $\{S_5=s\}$ for $s\in\{-5,-3,-1,1,3,5\}$. On it, $S_8$ has the conditional pmf $p_{S_8\mid S_5}(t\mid s)=p_{S_5,S_8}(s,t)/p_{S_5}(s)$, which needs the joint pmf of $(S_5,S_8)$: count paths.
    For $s=3$: $t=0,2,4,6$ with probabilities $\tfrac18,\tfrac38,\tfrac38,\tfrac18$, mean $3$, second moment $12$. Then five more rows.
  \end{block}
  \begin{block}{What conditioning offers}
    $S_8=S_5+D$ with $D=\xi_6+\xi_7+\xi_8$. On row $\{S_5=s\}$ the first term is the constant $s$ and comes out. The second is built from steps the row says nothing about, so the row should not change its distribution: $\E[D\mid S_5]$ should be $\E[D]=0$.
  \end{block}
  \begin{itemize}
    \item ``The row says nothing about $D$'' is $D\indep S_5$, which holds by the fact on the previous slide. The next slide says exactly what it buys.
  \end{itemize}
\end{frame}

% ---------------------------------------------------------------------------
\begin{frame}{Property 3: Independence. Every Row Looks Like the Whole Table}
  \small
  \begin{block}{Theorem}
    If $X\indep Y$ then $\E[X\mid Y]=\E[X]$, a constant; equivalently $\E[X\mid Y=y]=\E[X]$ for every $y\in S_Y$.
  \end{block}
  \begin{block}{Proof (discrete case)}
    Take $A=\{x\}$, $B=\{y\}$ in the definition: $p_{X,Y}(x,y)=p_X(x)\,p_Y(y)$. So for every $y\in S_Y$,
    \[
      p_{X\mid Y}(x\mid y)=\frac{p_X(x)\,p_Y(y)}{p_Y(y)}=p_X(x):
    \]
    each row carries the same distribution of $X$ as the whole table. Hence every row average is $\sum_x x\,p_X(x)=\E[X]$.
    Continuous case identical, with $f_{X,Y}=f_Xf_Y$. \hfill$\square$
  \end{block}
  \words{Knowing the row tells you nothing about $X$, so the row average is the overall average.}
  \begin{itemize}
    \item Two dice: $Z\indep Y$; every row of the $Z$-table is $1,\dots,6$, so $\E[Z\mid Y]=3.5$ on every row.
    \item Also $\E[h(X)\mid Y]=\E[h(X)]$ for $h\colon\R\to\R$, since $h(X)\indep Y$ too ($\{h(X)\in A\}=\{X\in h^{-1}(A)\}$).
  \end{itemize}
\end{frame}

% ---------------------------------------------------------------------------
\begin{frame}{Property 3: The Random Walk Done, and a Recipe}
  \small
  \label{rwdone}%
  \begin{exampleblock}{$\E[S_8\mid S_5]$ and $\E[S_8^2\mid S_5]$}
    $D=\xi_6+\xi_7+\xi_8$ takes the values $-3,-1,1,3$ with probabilities $\tfrac18,\tfrac38,\tfrac38,\tfrac18$, so $\E[D]=0$ and $\E[D^2]=9\cdot\tfrac14+1\cdot\tfrac34=3$; and $D\indep S_5$ (disjoint groups of steps).
    \begin{align*}
      \E[S_8\mid S_5]&=\E[S_5\mid S_5]+\E[D\mid S_5]=S_5+\E[D]=S_5,\\
      \E[S_8^2\mid S_5]&=\E[S_5^2+2S_5D+D^2\mid S_5]=S_5^2+2S_5\,\E[D]+\E[D^2]=S_5^2+3 .
    \end{align*}
    Check on row $s=3$: $\tfrac18\cdot0+\tfrac38\cdot4+\tfrac38\cdot16+\tfrac18\cdot36=12=3^2+3$.
  \end{exampleblock}
  \begin{block}{The recipe}
    To compute $\E[X\mid Y]$ without listing rows, write $X$ as sums and products of \emph{functions of $Y$} and \emph{quantities independent of $Y$}. The former come out unchanged (taking out what is known); the latter are replaced by their ordinary expectations (independence).
  \end{block}
  \begin{itemize}
    \item Two dice in one line: $X=Y+Z$ with $Z\indep Y$, so $\E[X\mid Y]=Y+\E[Z]=Y+3.5$; and $\E[X^2\mid Y]=Y^2+2Y\,\E[Z]+\E[Z^2]$.
    \item $\E[S_8\mid S_5]=S_5$: from wherever the walk stands, its expected future position is where it is. This is the martingale property, taken up in the stochastic-processes lectures.
  \end{itemize}
\end{frame}

% ---------------------------------------------------------------------------
\begin{frame}{Property 3: What Independence Does Not Say}
  \small
  \begin{block}{The converse fails: $\E[X\mid Y]=\E[X]$ does not force $X\indep Y$}
    $Y$ uniform on $\{-1,0,1\}$ and $W=Y^2$. Conditioning on $W$ creates two rows:
    \begin{center}
      \renewcommand{\arraystretch}{1.2}
      \begin{tabular}{@{}lccc@{}}
        \toprule
        row & outcomes & values of $Y$ & row average \\
        \midrule
        $\{W=0\}$ & $\{0\}$ & $0$ & $0$ \\
        $\{W=1\}$ & $\{-1,1\}$ & $-1,\ 1$ & $0$ \\
        \bottomrule
      \end{tabular}
    \end{center}
    So $\E[Y\mid W]=0=\E[Y]$, yet $W$ is a function of $Y$: $\Prob(Y=0,W=0)=\tfrac13\ne\Prob(Y=0)\Prob(W=0)=\tfrac19$.
  \end{block}
  \begin{itemize}
    \item Equal row \emph{averages} do not mean equal row \emph{distributions}. Independence is the second; the theorem only needs, and $\E[X\mid Y]=\E[X]$ only delivers, the first.
    \item What $\E[X\mid Y]=\E[X]$ does deliver, by tower and taking out: $\E[h(Y)X]=\E\big[h(Y)\,\E[X\mid Y]\big]=\E[h(Y)]\,\E[X]$, i.e.\ $\Cov(h(Y),X)=0$ for every $h\colon\R\to\R$. Here $\Cov(W,Y)=\E[Y^3]-\E[W]\E[Y]=0$.
  \end{itemize}
\end{frame}

% ---------------------------------------------------------------------------
\begin{frame}{Properties: Summary}
  \footnotesize
  \begin{center}
    \renewcommand{\arraystretch}{1.5}
    \begin{tabular}{@{}l>{$}l<{$}p{4.1cm}@{}}
      \toprule
      property & \text{statement} & one-glance reason \\
      \midrule
      tower & \E[\E[X\mid Y]]=\E[X] & painting rows with their averages keeps the total \\
      linearity & \E[aX+bX'\mid Y]=a\E[X\mid Y]+b\E[X'\mid Y] & averages are linear, row by row \\
      taking out what is known & \E[h(Y)X\mid Y]=h(Y)\E[X\mid Y] & $h(Y)$ is constant on a row \\
      known variables & \E[h(Y)\mid Y]=h(Y),\ \ \E[X\mid X]=X & a constant's average is itself \\
      independence & X\indep Y\Rightarrow\E[X\mid Y]=\E[X] & every row has the table's distribution of $X$ \\
      \bottomrule
    \end{tabular}
  \end{center}
  \begin{center}
    \renewcommand{\arraystretch}{1.4}
    \begin{tabular}{@{}lp{5.6cm}p{5cm}@{}}
      \toprule
      object & type & meaning \\
      \midrule
      $\E[X]$ & number & average over the whole table \\
      $\E[X\mid Y=y]$ & number for each $y\in S_Y$; a function $g\colon S_Y\to\R$ & average over row $y$ \\
      $\E[X\mid Y]$ & random variable $g(Y)\colon\Omega\to\R$ & row average, written into every cell of the row \\
      \bottomrule
    \end{tabular}
  \end{center}
\end{frame}

% ===========================================================================
\begin{frame}{Application 1: Random Sums. Problem B}
  \small
  Let $N\colon\Omega\to\{0,1,2,\dots\}$ and let $X_1,X_2,\dots\colon\Omega\to\R$ have $\E[X_i]=\mu$ for all $i$. Put
  \[
    S=X_1+X_2+\dots+X_N\qquad(S=0\text{ when }N=0).
  \]
  \words{Two layers of randomness: how many terms, and what each term is.}
  \begin{block}{Problem B, and why it is hard head-on}
    $N$ claims, $\E[N]=20$; claim $i$ has size $X_i$ with $\E[X_i]=\pounds1500$. Find $\E[S]$.
    The pmf or pdf of $S$ is $\sum_n\Prob(N=n)\cdot(\text{distribution of }X_1+\dots+X_n)$, an $n$-fold convolution for each $n$; and here neither the distribution of $N$ nor that of the $X_i$ is even specified. Only the two means are.
  \end{block}
  \begin{block}{What conditioning offers}
    On the row $\{N=n\}$, $S=X_1+\dots+X_n$ is a sum of a \emph{fixed} number of terms, and $\E[X_1+\dots+X_n]=n\mu$ by linearity. So $\E[S\mid N=n]$ would be $n\mu$, provided that restricting to the row does not change the distribution of the terms.
  \end{block}
  \begin{itemize}
    \item That proviso is a genuine assumption about how $N$ is decided. It is stated next.
  \end{itemize}
\end{frame}

% ---------------------------------------------------------------------------
\begin{frame}{Application 1: Random Sums. The Assumption}
  \small
  \begin{block}{Assumption: $N$ is independent of the $X_i$}
    For every $n$ and all sets $A_1,\dots,A_n\subseteq\R$,
    \[
      \Prob(N=n,\ X_1\in A_1,\dots,X_n\in A_n)=\Prob(N=n)\,\Prob(X_1\in A_1,\dots,X_n\in A_n).
    \]
  \end{block}
  \words{Given $N=n$, the joint distribution of $(X_1,\dots,X_n)$ is what it was unconditionally.}
  \begin{itemize}
    \item Not assumed: independence of the $X_i$ from each other, or identical distributions. Only equal means $\mu$.
    \item Holds when $N$ is settled by a separate mechanism: the number of customers before anyone spends; the number of claims regardless of their sizes.
    \item Fails when $N$ is decided by looking at the $X_i$: ``stop at the first six''; ``include the second term only if it is small'' (slide~\ref{insurer}).
  \end{itemize}
\end{frame}

% ---------------------------------------------------------------------------
\begin{frame}{Application 1: Random Sums. Theorem and Proof}
  \small
  \begin{block}{Theorem}
    Under the assumption, $\E[S\mid N=n]=n\mu$ for every $n$, hence $\E[S\mid N]=\mu N$ and $\E[S]=\mu\,\E[N]$.
  \end{block}
  \begin{block}{Proof}
    On the row $\{N=n\}$, $S=W_n:=X_1+\dots+X_n$. By the assumption, the row $\{N=n\}$ does not change the distribution of $W_n$, so Property 3, applied to $W_n$ at the row $n$, gives
    \[
      \E[S\mid N=n]=\E[W_n\mid N=n]=\E[W_n]=n\mu
    \]
    by linearity. Substitute $n=N$ to get $\E[S\mid N]=\mu N$; then $\E[S]=\E[\mu N]=\mu\,\E[N]$ by the tower property. \hfill$\square$
  \end{block}
  \words{Given the row, the random number of terms is a fixed number and linearity does the rest.}
  \begin{itemize}
    \item The three steps again: condition on the layer that makes the problem hard; on the row it is gone; tower.
    \item Nothing about the distribution of $S$ was computed. Its mean is all that was asked for and all that was found.
  \end{itemize}
\end{frame}

% ---------------------------------------------------------------------------
\begin{frame}{Application 1: Back to the Insurer, and When the Assumption Fails}
  \small
  \label{insurer}%
  \begin{exampleblock}{Problem B}
    $\E[S]=\mu\,\E[N]=1500\cdot20=\pounds30\,000$, whatever the distributions of $N$ and of the claim sizes, as long as the count is independent of the sizes.
  \end{exampleblock}
  \begin{exampleblock}{$N$ decided by peeking at a term: the formula gives the wrong number}
    Two dice $X_1,X_2$. Let $N=1$ if $X_2=6$ and $N=2$ otherwise (the second die counts only if it is not a six). Then $S=X_1+X_2\,\ind_{\{X_2\ne6\}}$ and
    \[
      \E[S]=3.5+\tfrac{1+2+3+4+5}6=6,\qquad\text{but}\qquad \mu\,\E[N]=3.5\cdot\tfrac{11}6=\tfrac{77}{12}\approx6.42 .
    \]
    On the row $\{N=2\}=\{X_2\ne6\}$ the second term averages $3$, not $\mu=3.5$: the row changed the distribution of a term.
  \end{exampleblock}
  \begin{itemize}
    \item ``Stop at the first six'' also violates the assumption, yet $\mu\,\E[N]=3.5\cdot6=21$ is right. That is Wald's identity, a separate theorem: its hypothesis is that $N$ never looks \emph{ahead}. The peeking example looks ahead.
  \end{itemize}
\end{frame}

% ===========================================================================
\begin{frame}{Application 2: Waiting for HH. Problem C}
  \small
  Toss a fair coin repeatedly. Let $T\colon\Omega\to\{2,3,\dots\}$ be the number of the toss on which two consecutive heads are first completed: for T\,H\,T\,H\,H\,\dots, $T=5$. We want $\E[T]$.
  \begin{block}{Why head-on is hard}
    $\Prob(T=n)$ is the number of H/T strings of length $n$ that end in HH and contain no earlier HH, over $2^n$. The count is the Fibonacci number $F_{n-1}$, so $\Prob(T=n)=F_{n-1}/2^n$, and $\E[T]=\sum_{n\ge2}nF_{n-1}/2^n$ needs generating functions to sum.
  \end{block}
  The rows we condition on here are described by events rather than by the value of a named variable, so:
  \begin{block}{Definition (conditional expectation given an event)}
    For $A\in\F$ with $\Prob(A)>0$ and $\ind_A\colon\Omega\to\{0,1\}$ the indicator of $A$,
    \[
      \E[X\mid A]=\sum_x x\,\Prob(X=x\mid A)=\E[X\mid \ind_A=1].
    \]
  \end{block}
  \words{A row can be any event of positive probability: restrict to $A$, rescale, average.}
\end{frame}

% ---------------------------------------------------------------------------
\begin{frame}{Application 2: Waiting for HH. The Restart Property}
  \small
  Let $A_k$ be the event that the first $k$ tosses form a given H/T string ending in T; $\Prob(A_k)=2^{-k}>0$. On $A_k$ no progress towards HH has been made at toss $k$.
  \begin{block}{Restart property}
    Let $T'$ be the number of tosses after toss $k$ until HH is first completed among tosses $k+1,k+2,\dots$. Then on $A_k$, $T=k+T'$, and
    \[
      \E[T\mid A_k]=k+\E[T'\mid A_k]=k+\E[T].
    \]
  \end{block}
  \begin{block}{Why}
    $T'$ is a function of tosses $k+1,k+2,\dots$, and $\ind_{A_k}$ of tosses $1,\dots,k$: disjoint groups of independent tosses, so $T'\indep\ind_{A_k}$ and, by Property 3, $\E[T'\mid A_k]=\E[T']$. The tosses $k+1,k+2,\dots$ are a fair coin tossed repeatedly, exactly the experiment defining $T$, so $T'$ has the distribution of $T$ and $\E[T']=\E[T]$. \hfill$\square$
  \end{block}
  \words{On such a row $T$ is $k$ plus a fresh copy of $T$, so the row average is $k$ plus the very unknown we are after.}
\end{frame}

% ---------------------------------------------------------------------------
\begin{frame}{Application 2: Waiting for HH. The Computation}
  \small
  \label{hh}%
  \begin{columns}[T]
    \begin{column}{0.6\textwidth}
      \textbf{The conditioning variable.} Let $Y\colon\Omega\to\{1,2,3\}$ record how the game opens:
      \[
        Y=\begin{cases}1&\text{toss 1 is T}\\2&\text{tosses 1,2 are H,T}\\3&\text{tosses 1,2 are H,H}\end{cases}
      \]
      with $p_Y(1)=\tfrac12$, $p_Y(2)=p_Y(3)=\tfrac14$.
      Every outcome falls in exactly one case: three rows.
      \medskip

      \textbf{Row averages.} Rows 1 and 2 are the events $A_1=\{\text{T}\}$ and $A_2=\{\text{HT}\}$ (restart); row 3 is finished:
      \[
        \E[T\mid Y=1]=1+\E[T],\qquad \E[T\mid Y=2]=2+\E[T],
      \]
      and $\E[T\mid Y=3]=2$.
      \textbf{Total expectation.}
      \begin{align*}
        \E[T]&=\tfrac12\big(1+\E[T]\big)+\tfrac14\big(2+\E[T]\big)+\tfrac14\cdot2\\
             &=\tfrac34\E[T]+\tfrac32,\qquad\text{so }\E[T]=6 .
      \end{align*}
    \end{column}
    \begin{column}{0.38\textwidth}
      \centering
      \begin{tikzpicture}[font=\scriptsize,
          box/.style={draw,rounded corners,inner sep=2pt,fill=white},
          lbl/.style={inner sep=1.5pt,fill=white,midway}]
        \node[box] (s) at (0,0) {start};
        \node[box,fill=shade] (t) at (-1.2,-1.3) {T: restart};
        \node[box] (h) at (1.2,-1.3) {H};
        \node[box,fill=shade] (ht) at (0.2,-2.6) {HT: restart};
        \node[box,fill=leaf!20] (hh) at (2.2,-2.6) {HH: done};
        \draw (s) -- node[lbl] {$\tfrac12$} (t);
        \draw (s) -- node[lbl] {$\tfrac12$} (h);
        \draw (h) -- node[lbl] {$\tfrac12$} (ht);
        \draw (h) -- node[lbl] {$\tfrac12$} (hh);
      \end{tikzpicture}
      \par\medskip
      {\footnotesize Two of the three leaves lead back to the start. That is why $\E[T]$ appears on both sides.}
    \end{column}
  \end{columns}
\end{frame}

% ---------------------------------------------------------------------------
\begin{frame}{Application 2: The Same Move, Elsewhere}
  \small
  Whenever a process restarts after certain outcomes, condition on the first step: the unknown mean appears on both sides of the total-expectation equation, and the equation solves it.
  \begin{itemize}
    \item \textbf{Rolls until the first six}, $T_6$: \ $\E[T_6]=\tfrac16\cdot1+\tfrac56\big(1+\E[T_6]\big)$, so $\E[T_6]=6$.
    \item \textbf{Waiting for a T}, $T_{\mathrm T}$: the same equation with $\tfrac12$ for $\tfrac16$ gives $\E[T_{\mathrm T}]=2$.
    \item \textbf{Pattern HT}, $T_{\mathrm{HT}}$: after the first H, wait for a T. \ $\E[T_{\mathrm{HT}}]=\tfrac12\big(1+\E[T_{\mathrm{HT}}]\big)+\tfrac12\big(1+\E[T_{\mathrm T}]\big)$, so $\E[T_{\mathrm{HT}}]=4<6$: when HT fails (HH), the second H is a head start; when HH fails (HT), nothing is.
    \item \textbf{Pattern HHH:} $14$. \ \textbf{Two sixes in a row:} $42$. Same equation, more rows.
    \item \textbf{No memory:} $\E[T_6\mid T_6>3]=3+\E[T_6]=9$, by the restart property at $k=3$.
  \end{itemize}
  \begin{block}{The equation is legitimate}
    Solving $\E[T]=\tfrac34\E[T]+\tfrac32$ presumes $\E[T]<\infty$. Each block of two tosses is HH with probability $\tfrac14$, so $\Prob(T>2m)\le(\tfrac34)^m$, and $\E[T]=\sum_{n\ge0}\Prob(T>n)<\infty$ (tail-sum formula for $\{0,1,2,\dots\}$-valued $T$; standard, not proved here).
  \end{block}
\end{frame}

